\documentclass[11pt]{article}
\usepackage{cosai}

\title{Add CoSAI Full Paper Title Here}
\author{Workstream Name}

\begin{document}

% --- Front Cover ---
\makeCoSAICover{ADD TITLE HERE}{Workstream Name}

% --- Table of Contents ---
\newpage
\tableofcontents

% --- Abstract Page ---
\newpage
\section*{Abstract}
\begin{center}
    \begin{tabular}{ccc}
        \sectionIcon{a}{Practical}{Cooperation} & 
        \sectionIcon{b}{MCP}{Security} & 
        \sectionIcon{}{Actionable}{Controls}
    \end{tabular}
\end{center}

\vspace{1cm}
Summary of the technical purpose of the document. [cite_start]By creating structured opportunities for interaction and feedback, it aims to foster cooperation and produce outcomes that are practical, measurable, and valued by all involved. [cite: 139]

% --- Introduction (Dark Background Section) ---
\newpage
\pagecolor{cosaiDarkBlue}\color{white}
\section*{\color{white}Introduction}
[cite_start]Since its emergence a year ago, MCP has rapidly established itself as the protocol for transmitting structured context between AI agents and services. [cite: 147] [cite_start]Given the growing importance and attack surface of MCP and agentic systems, it is imperative that deployment specific security threats are identified. [cite: 148]

% --- Technical Section ---
\newpage
\pagecolor{white}\color{black}
\section{MCP Architecture and Standards}
[cite_start]MCP is an open standard developed by Anthropic and a growing open source community that provides a structured framework for connecting large language models and AI agents to external tools, data sources, and services. [cite: 154]

\subsection{Client-Server Architecture}
[cite_start]MCP follows a client-server architecture where host applications (such as AI assistants, IDEs, or workflow automation tools) use MCP clients to connect to local or remote MCP servers. [cite: 166]

\begin{itemize}
    [cite_start]\item \textbf{Discovery:} Standardized methods for discovering available capabilities. [cite: 159]
    [cite_start]\item \textbf{Invocation:} Invoking parametrized services. [cite: 160]
    [cite_start]\item \textbf{Transport:} Supports stdio for local and Streamable HTTP for remote. [cite: 169]
\end{itemize}

% --- Conclusions ---
\section{Takeaways and Conclusion}
[cite_start]Organizations deploying MCP-based systems must develop defense-in-depth strategies including zero-trust architectures and hardware-based isolation through trusted execution environments. [cite: 190]

% --- Acknowledgements ---
\newpage
\section*{Acknowledgements}
\textbf{Workstream Leads}
\begin{itemize}
    [cite_start]\item Sarah Novotny [cite: 199]
    [cite_start]\item Ian Molloy, IBM [cite: 200]
    [cite_start]\item Raghu Yeluri, Intel [cite: 201]
\end{itemize}

\end{document}